% Generated by Listas/make_statements.py from pratica-lecture-04.md, item 4. Do not edit.
\textbf{Source.} Muñoz-Garcia \& Toro-Gonzalez (2019), \emph{Strategy and Game Theory: Practice Exercises with Answers}, Ch. 4, Exercise 11 ``Stackelberg with Two Firms'', p.~180 {[}PDF 192{]}. Reproduced verbatim (statement only; the book's worked answer follows it).



\begin{quote}

Consider a leader and a follower in a Stackelberg game of quantity competition. Firms face an inverse demand curve \(p(Q) = 1 - Q\), where \(Q = q_L + q_F\) denotes aggregate output. The leader faces a constant marginal cost \(c_L > 0\) while the follower's constant marginal cost is \(c_F > 0\), where \(1 > c_F \geq c_L\), indicating that the leader has access to a more efficient technology than the follower.



\textbf{Part (a)} Find the follower's best response function, \(BRF_F\), i.e., \(q_F(q_L)\).



\textbf{Part (b)} Determine each firm's output strategy in the subgame perfect equilibrium (SPNE) of this sequential-move game.



\textbf{Part (c)} Under which conditions on \(c_L\) can you guarantee that both firms produce strictly positive output levels?



\textbf{Part (d)} Assuming cost symmetry, i.e., \(c_L = c_F = c\), determine the aggregate equilibrium output level in the Stackelberg game, and its associated equilibrium price. Then compare them with those arising in (1) a monopoly; in (2) the Cournot game of quantity competition; and (3) in the Bertrand game of price competition (which, under cost symmetry, is analogous to a perfectly competitive industry) {[}\emph{Hint}: For a more direct and visual comparison, depict the inverse demand curve \(p(Q) = 1 - Q\), and locate each of these four aggregate output levels in the horizontal axis (with its associated prices in the vertical axis).{]}.

\end{quote}



\begin{center}\rule{0.5\linewidth}{0.5pt}\end{center}

